\ProvidesFile{ch-front.tex}[2024-09-12 front matter chapter]
%
%  This is the ``front matter'' for the thesis.
%
%  REFERENCES
%
%    TCMOS17
%      The Chicago Manual of Style Online, 17th edition.
%      https://www.chicagomanualofstyle.org/home.html
%      retrieved on 2020-02-29
%
%    TEMPL
%      Thesis and Disertation Office Templates.
%      https://www.purdue.edu/gradschool/research/thesis/templates.html
%      retrieved on 2020-02-29
%
%    WNNCD
%    Webster's Ninth New Collegiate Dictionary.
%

%
%   Only Purdue University uses this page
%
%   Comment out \begin{statement} through \end{statement}
%   if you are not at Purdue University.
%
% Statement of Thesis/Dissertation Approval Page
% This page is REQUIRED.  The page should be numbered "2"
% and should NOT be listed in your TABLE OF CONTENTS.
\begin{statement}
  % Delete or add \entry commands as needed for all committe members.
  \entry{Dr.~Christina Garman, Chair}{Department of Computer Science}
  \entry{Dr.~Jeremiah Blocki}{Department of Computer Science}
  \entry{Dr.~Jing (Dave) Tian}{Department of Computer Science}
  \entry{Dr.~Antonio Bianchi}{Department of Computer Science}
  % There should be one \approvedby command containing the
  % "FORM 9 THESIS FORM HEAD NAME HERE" (from TEMPL, retrieved on 2020-03-01).
  \approvedby{Dr.~Voicu Popescu}
\end{statement}

% Dedication page is optional.
% A name and often a message in tribute to a person or cause.
% References: WEB9 332.
% \begin{dedication}
%   To graduate students
% \end{dedication}

% Acknowledgements page is optional but most theses include
% a brief statement of appreciation or recognition of special
% assistance.
\begin{acknowledgments}
First and foremost, I owe my deepest gratitude to my doctoral advisor and mentor, Dr. Christina Garman. The opportunity to pursue my research under her guidance has shaped not only this dissertation, but also the researcher and educator I have become. She gave me the opportunity to explore questions that genuinely interested me, while always helping me recognize which ideas were worth pursuing and how they fit into a larger research vision. Her mentorship taught me to approach difficult problems with independence, patience, and care. Just as importantly, her guidance extended far beyond individual projects or academic milestones. Through our conversations, I came to understand my own goals more clearly, develop confidence in my abilities, and reflect on the kind of academic career I hope to build. I feel extraordinarily fortunate to have had an advisor who supported my growth with such generosity, trust, patience, and encouragement.

I am also sincerely grateful to Dr. Jeremiah Blocki, Dr. Dave (Jing) Tian, and Dr. Antonio Bianchi for serving on my dissertation committee. Their thoughtful questions and careful feedback continually pushed me to examine my assumptions, clarify my arguments, and view my work from broader perspectives. Their guidance strengthened this dissertation and, more importantly, deepened the way I think about research.

My path toward this Ph.D. began long before I arrived at Purdue, and it was shaped by many faculty members who offered me opportunities, guidance, and encouragement at different stages of my academic journey. I am grateful to Dr. Xiaojing Liao of the University of Illinois Urbana-Champaign, Dr. David Crandall, Dr. Selma Sabanovic, and Dr. William Swanson of Indiana University Bloomington, and Dr. James Elder of York University. Each of them influenced the way I approach research and helped establish the foundation on which my later work was built. I also deeply appreciate Dr. Jennifer Coy of Ball State University and Dr. Yue Xiao of William \& Mary for their advice and support throughout my academic career. I am especially grateful to Dr. Xiao, whose assistance and encouragement have meant a great deal to me during important moments along this path.

Working alongside generous and talented collaborators was one of the most rewarding aspects of my doctoral experience. I am deeply grateful to Dr. Priyam Biswas of Intel, Dr. Rui Zhu of Amazon, and Yuquan Xu of Georgia Tech. Our collaborations were built through shared ideas, long discussions, technical challenges, and collective effort. I learned a great deal from each of you, and I sincerely value both the work we accomplished together and the experience of learning from one another.

I was also fortunate to be part of the Boilermakers' Applied Research in Cryptography (BARC) group. To Alex Seto, Jacob White, Arushi Arora, and Saurav Chittal, thank you for the questions, feedback, conversations, and companionship that filled our weekly meetings and many discussions beyond them. Research can often be uncertain and isolating, but sharing ideas with a group of thoughtful and supportive colleagues made the process both more productive and more enjoyable. I will remember not only the technical insights we exchanged, but also the sense of community we created together.

Finally, I want to thank my friends Chuck Jia, Sijing Peng, Yiwei Zhang, and Chenchen Zheng who stood beside me through moments of uncertainty, exhaustion, change, and celebration. I am also deeply grateful to Pankil Bhatt, Dr. Winona Snapp-Childs, Wayman Mai, Charles Pope, and Dr. Olgun Sadik, whose guidance, encouragement, and support helped me navigate important academic and professional decisions throughout this journey. I sincerely appreciate the friendship, mentorship, generosity, and care I received from all of them. Their presence brought comfort, perspective, and joy beyond my academic work, while their advice and support helped me grow both personally and professionally. This journey would have been far more difficult, and far less meaningful, without them.

This work may bear my name, but the journey that produced it was never mine alone. It was shaped by the faculties who opened doors for me, the mentors who believed in me, the collaborators and colleagues who challenged and inspired me, and the friends who gave me strength when I needed it most. To all of them, I offer my heartfelt gratitude.
\end{acknowledgments}

% The preface is optional.
% References: TCMOS17 1.49, WEB9 927.
% \begin{preface}
%   This is the preface.
% \end{preface}

% The Table of Contents is required.
% The Table of Contents will be automatically created for you
% using information you supply in
%     \chapter
%     \section
%     \subsection
%     \subsubsection
%     commands.
\pdfbookmark{TABLE OF CONTENTS}{Contents}
\tableofcontents

% If your thesis has tables, a list of tables is required.
% The List of Tables will be automatically created for you using
% information you supply in
%     \begin{table} ... \end{table}
% environments.
\listoftables

% If your thesis has figures, a list of figures is required.
% The List of Figures will be automatically created for you using
% information you supply in
%     \begin{figure} ... \end{figure}
% environments.
\listoffigures

% If your thesis has listings, a list of listings is required.
% The List of Listings will be automatically created for you using
% information you supply in
%     \begin{ZZlisting} ... \end{ZZlisting}
% environments.
\ZZlistoflistings

% If your thesis has protocols, you may want to do a list of protocols.
% The List of Protocols will be automatically created for you using
% information you supply in
%     \begin{protocol} ... \end{protocol}
% environments.
\listofprotocols

% If your thesis has schemes, you may want to do a list of schemes.
% The List of Schemes will be automatically created for you using
% information you supply in
%     \begin{scheme} ... \end{scheme}
% environments.
% \listofschemes

% List of Symbols is optional.
% \begin{symbols}
%   \(m\)& mass\cr
%   \(v\)& velocity\cr
% \end{symbols}

% List of Abbreviations is optional.
% \begin{abbreviations}
%   abbr& abbreviation\cr
%   bcf& billion cubic feet\cr
%   BMOC& big man on campus\cr
% \end{abbreviations}

% Nomenclature is optional.
% \begin{nomenclature}
%   alanine& 2-Aminopropanoic acid\\
%   gasoline& a transparent,
%     petroleum-derived flammable liquid that is used primarily
%     as a fuel in most spark-ignited internal combustion engines\\
%   valine& 2-Amino-3-methylbutanoic acid\\
%   Valvoline& Valvoline Inc.~is an American manufacturer
%     and distributor of Valvoline-brand automotive oil,
%     additives,
%     and lubricants.
%     It also owns the Valvoline Instant Oil Change
%     and Valvoline Express Care chains of car repair centers.
%     As of 2016,
%     it is the second largest oil change service provider in the United States
%     with 10\% market share
%     and 1,050 locations.\\
%   \mbox{}\\
%   \rlap{you can divide these into categories if you want}\\
%   \mbox{}\\
%   \rlap{\bfseries Biology}\\
%   alanine& 2-Aminopropanoic acid\\
%   valine& 2-Amino-3-methylbutanoic acid\\
%   \rlap{\bfseries Transportation}\\
%   gasoline& a transparent,
%     petroleum-derived flammable liquid that is used primarily
%     as a fuel in most spark-ignited internal combustion engines\\
%   Valvoline& Valvoline Inc.~is an American manufacturer
%     and distributor of Valvoline-brand automotive oil,
%     additives,
%     and lubricants.
%     It also owns the Valvoline Instant Oil Change
%     and Valvoline Express Care chains of car repair centers.
%     As of 2016,
%     it is the second largest oil change service provider in the United States
%     with 10\% market share
%     and 1,050 locations.\\
% \end{nomenclature}

% Glossary is optional.
\begin{glossary}
  Binary Analysis & Binary analysis is the examination of compiled executables without source code to understand program structure, behavior, and embedded functionality at the machine instruction level.
\\
  Black-box Fuzzing & A testing approach where inputs are generated without any knowledge of the program's internal structure. The system is treated as opaque, and only external behaviors such as crashes or errors are observed.
\\
  Cryptographic Logic Error & An error or bug that causes a program to produce an incorrect result due to errors in the cryptographic protocol implementation, such as incorrect mathematical operations or disregarding the specification.
\\
  Cryptographic Primitive & Cryptographic primitive (also called cryptographic function or cryptographic algorithm) refers to the entirety of a specific scheme or class of algorithms such as AES or RSA, used to generate ciphertext and support encryption and decryption processes for securing information.
\\
  Differential Fuzzing & A software testing technique that detects bugs by applying identical inputs to comparable programs or separate implementations and examining differences in their behavior.
\\
  Grey-box Fuzzing & A testing approach that uses partial knowledge of the program's internal behavior. Lightweight runtime feedback, such as code coverage or execution paths, is used without requiring full program analysis or source-code access.
\\
  Taint Tracking & A program-analysis technique that marks selected input data as tainted and tracks how that data propagates through instructions, variables, memory locations, and control flow.
\\
  White-box Fuzzing & A testing approach where inputs are generated using full knowledge of the program's internal structure. Code and execution paths are analyzed to systematically explore behaviors and improve test coverage.
\\
  Zero-Knowledge Proof & A cryptographic protocol that allows a user (the prover) proves a statement to another party (the verifier) without revealing anything about the statement other than that it is true.
\\
  zk-SNARKs & Zero-Knowledge Succinct Non-Interactive Arguments of Knowledge are a non-interactive zero-knowledge proof system with a short proof size and efficient verification, enabling practical deployment in cryptographic systems.

\end{glossary}

% Abstract is required.
% Reference:
%       author    = {Thesis and Dissertation Office},
%       pages     = {9},
%       publisher = {Purdue University},
%       url       = {https://www.purdue.edu/gradschool/research/thesis/templates.html},
%       urldate   = {2022-10-24},
% Choose:
%     Microsoft Word
%     West Lafayette
%     College of Engineering

\begin{abstract}

Cryptographic software forms a critical foundation of modern computing systems, but the security guarantees of cryptographic protocols do not automatically extend to their implementations. Errors in arithmetic operations, validation logic, data conversion, constraint generation, or component integration can cause deployed software to deviate from the intended protocol while still producing plausible outputs. Such risks are difficult to detect in compiled binaries and become even more challenging in modern cryptographic systems such as zero-knowledge proofs, where implementations combine finite-field arithmetic, constraint systems, witness generation, proving procedures, verification logic, and serialization formats.

Securing cryptographic implementations requires analysis techniques that can reason about both low-level program behavior and high-level cryptographic intent. To address this need, this dissertation first examines cryptographic function identification in binaries. It categorizes existing detection techniques, develops a unified benchmarking framework, and evaluates current tools through reproduction and replication studies across different compilers, optimization levels, obfuscation strategies, and algorithm variants. The second part introduces an automated security analysis framework for zkSNARK implementations that combines constraint checking with fuzzing-based testing to detect and locate cryptographic logic errors. This approach helps determine whether an implemented constraint system correctly enforces the intended computation. The third part develops a black-box and grey-box differential fuzzing approach for zero-knowledge proof binary applications. It uses structured input generation, coverage monitoring, control-dependency-aware taint tracking, and error localization to guide testing toward security-relevant code and expose inconsistencies in circuit validation, witness conversion, public-input handling, verification logic, and final proof checking.

Together, these contributions connect binary analysis, automated checking, and protocol-aware fuzzing to improve the practical security of cryptographic software. They provide methods for identifying implementation-level weaknesses that may remain hidden during ordinary testing, especially when programs produce valid-looking outputs despite incorrect cryptographic behavior. By combining systematic evaluation, zkSNARK-specific analysis, and binary-level testing, this dissertation supports the development of more reliable techniques for analyzing, testing, and securing real-world cryptographic systems.

\end{abstract}
